\par\Needspace{6\baselineskip}\medskip
\hypertarget{native-081-article062}{}\pdfbookmark[2]{62.}{article-062}
\begin{center}62.\end{center}
Sit corporis coelestis longitudo heliocentrica $\lambda$, latitudo $\beta$; longitudo geocentrica $l$, latitudo $b$, distantia a Sole $r$, a terra $\Delta$; denique terrae longitudo heliocentrica $L$, latitudo $B$, distantia a Sole $R$. Quum non statuamus $B=0$, formulae nostrae ad eum quoque casum applicari poterunt, ubi loci heliocentrici et geocentricus non ad eclipticam sed ad quodvis aliud planum referuntur, modo denominationes longitudinis et latitudinis supprimere oportebit: praeterea parallaxeos ratio statim haberi potest, si modo locus heliocentricus terrae non ad centrum sed ad locum in superficie immediate refertur. Statuamus porro $r\cos\beta=r'$, $\Delta\cos b=\Delta'$, $R\cos B=R'$. Iam referendo locum corporis coelestis atque terrae in spatio ad tria plana, quorum unum sit ecliptica, secundumque et tertium polos suos habeant in longitudine $N$ et $N+90^0$, protinus emergent aequationes sequentes:
\[
\begin{gathered}
r'\cos(\lambda-N)-R'\cos(L-N)=\Delta'\cos(l-N)\\
r'\sin(\lambda-N)-R'\sin(L-N)=\Delta'\sin(l-N)\\
r'\operatorname{tang}\beta-R'\operatorname{tang}B=\Delta'\operatorname{tang}b,
\end{gathered}
\]
ubi angulus $N$ omnino arbitrarius est. Aequatio prima et secunda statim determinabunt $l-N$ atque $\Delta'$, unde et ex tertia demanabit $b$; ex $b$ et $\Delta'$ habebis $\Delta$. Iam ut labor calculi quam commodissimus evadat, angulum arbitrarium $N$ tribus modis sequentibus determinamus:

I. Statuendo $N=L$, faciemus $\frac{r'}{R'}\sin(\lambda-L)=P$, $\frac{r'}{R'}\cos(\lambda-L)-1=Q$, invenienturque $l-L$, $\frac{\Delta'}{R'}$ atque $b$ per formulas
\[
\begin{gathered}
\operatorname{tang}(l-L)=\frac{P}{Q}\\[.5em]
\frac{\Delta'}{R'}=\frac{P}{\sin(l-L)}=\frac{Q}{\cos(l-L)}\\[.5em]
\operatorname{tang}b=\frac{\frac{r'}{R'}\operatorname{tang}\beta-\operatorname{tang}B}{\frac{\Delta'}{R'}}\mathord{\cdot}
\end{gathered}
\]

\SourceLeaf{082}II. Statuendo $N=\lambda$, faciemus $\frac{R'}{r'}\sin(\lambda-L)=P$, $1-\frac{R'}{r'}\cos(\lambda-L)=Q$, eritque
\[
\begin{gathered}
\operatorname{tang}(l-\lambda)=\frac{P}{Q}\\[.5em]
\frac{\Delta'}{r'}=\frac{P}{\sin(l-\lambda)}=\frac{Q}{\cos(l-\lambda)}\\[.5em]
\operatorname{tang}b=\frac{\operatorname{tang}\beta-\frac{R'}{r'}\operatorname{tang}B}{\frac{\Delta'}{r'}}\mathord{\cdot}
\end{gathered}
\]

III. Statuendo $N=\frac12(\lambda+L)$, invenientur $l$ atque $\Delta'$ per aequationes
\[
\begin{gathered}
\operatorname{tang}\{l-\tfrac12(\lambda+L)\}=\frac{r'+R'}{r'-R'}\operatorname{tang}\tfrac12(\lambda-L)\\[.5em]
\Delta'=\frac{(r'+R')\sin\frac12(\lambda-L)}{\sin(l-\frac12(\lambda+L))}=\frac{(r'-R')\cos\frac12(\lambda-L)}{\cos(l-\frac12(\lambda+L))}
\end{gathered}
\]
ac dein $b$ per aequationem supra datam. Logarithmus fractionis $\frac{r'+R'}{r'-R'}$ commode calculatur, si statuitur $\frac{R'}{r'}=\operatorname{tang}\zeta$, unde fit $\frac{r'+R'}{r'-R'}=\operatorname{tang}(45^0+\zeta)$. Hoc modo methodus III ad determinationem ipsius $l$ aliquanto brevior est, quam I et II, ad operationes reliquas autem has illi praeferendas censemus.
